% This is a conglomerate of various tex files I (Jan-Matthias Braun) merged
% for my thesis with excerpts from Jeremie Papon's Thesis (https://github.com/jpapon/papon_thesis)
% which itself was forked of the original Harvard Latex template by Jordan Suchow (suchow@fas.harvard.edu).
%
% In case of questions, please contact Jan-Matthias Braun (jan_braun@gmx.net).
% Otherwise feel free to use this template for your own thesis.
%
% Notes
% (1) This file is intended for use with xelatex, but the corresponding
%     lines for pdflatex are included as comments.
% (2) Please consult the comments on Layout-Tuning below.
% (3) This is the future, and it uses UTF-8, you can directly include
%     characters like ℝ, with the right definitions (see UTF-8 below).

%%%%%%%%%%%%%%%%%%%%%%%%%%%%%%%%%%%%%%%%%%%%%%%%%%%%%%%%%%%%%%%%%   Authors
%
% Original Harvard Template by Jordan Suchow (suchow@fas.harvard.edu)
%
% Fork by Jeremie Papon https://github.com/jpapon/papon_thesis
% (c) 2014
%
% Fork by Jan-Matthias Braun jan_braun@gmx.net
% (c) 2015
%
% Fork by Simon Reich https://github.com/simonreich
% (c) 2017

%%%%%%%%%%%%%%%%%%%%%%%%%%%%%%%%%%%%%%%%%%%%%%%%%%%%%%%%%%%%%%%%%   General Packages

\usepackage{calc}                       %% ermoeglicht Rechnen mit Laengen und Zaehlern
\usepackage{microtype}                  %% defaults are protrusion=true,expansion=true
\usepackage{quotchap}                   %% Quotation at chapters
\usepackage{setspace}                   %% line-spacing
\onehalfspace                           %% \onehalfspace corresponds to \setstretch{1.2}
% \doublespace                          %% To have more space for corrections.


% sidewaysfigure and rotations
\usepackage[figuresright]{rotating}

% Change the style of pages with huge floats with \thisfloatpagestyle
% http://tex.stackexchange.com/questions/169908/how-to-set-the-pagestyle-on-the-page-a-particular-float-ends-up-on
\usepackage{floatpag}
\usepackage{pdflscape}                  %% Set landscape mode for singular pages
\usepackage[authoryear,comma,sort&compress]{natbib}  %% Plain bibliography style
\bibliographystyle{plainnat}
\usepackage{units}                      %% units package



%%%%%%%%%%%%%%%%%%%%%%%%%%%%%%%%%%%%%%%%%%%%%%%%%%%%%%%%%%%%%%%%%%   For nice looking tables
\usepackage{booktabs}		            %% Befehle fuer besseres Tabellenlayout
\usepackage{longtable}		            %% umbrechbare Tabellen
\usepackage{multirow}



%%%%%%%%%%%%%%%%%%%%%%%%%%%%%%%%%%%%%%%%%%%%%%%%%%%%%%%%%%%%%%%%%%   tikz
\usepackage{tikz}
\usetikzlibrary{shapes,arrows,fit}
\usepackage{gnuplot-lua-tikz}



%%%%%%%%%%%%%%%%%%%%%%%%%%%%%%%%%%%%%%%%%%%%%%%%%%%%%%%%%%%%%%%%%%   Abbreviations
% Handle abbreviations properly
% This is useful for inline enumerations, e.g., an explicit 1. ...
\usepackage{xpunctuate}                 %% Define shortcuts for often used abbreviations
\newcommand{\ie}{i\xperiod{}e\xperiod{}}
\newcommand{\eg}{e\xperiod{}g\xperiod{}}
\newcommand{\etc}{etc\xperiod{}}
\newcommand{\ff}{ff\xperiod{}}

% Figure References
\newcommand{\FigurenameC}{Fig.}
\newcommand{\FigurenameS}{Fig.}
\newcommand{\Figref}[1]{\FigurenameC~\ref{#1}}
\newcommand{\figref}[1]{\FigurenameS~\ref{#1}}
\newcommand{\Tabref}[1]{Tab.~\ref{#1}}
\newcommand{\tabref}[1]{Tab.~\ref{#1}}



%%%%%%%%%%%%%%%%%%%%%%%%%%%%%%%%%%%%%%%%%%%%%%%%%%%%%%%%%%%%%%%%%   Font selection
% See https://wiki.ubuntuusers.de/XeTeX/
\usepackage{xltxtra}


% See http://www.macfreek.nl/memory/Fonts_in_LaTeX
% I have the impression, that Palatino is better to read on paper ... 
% On screen I prefer something else.
\setmainfont[Mapping=tex-text, Numbers=OldStyle]{TeX Gyre Pagella}      % Palatino as main font
%\setmainfont[Numbers=OldStyle]{TeX Gyre Termes}        % Times New Roman
\setsansfont[Mapping=tex-text, Numbers=OldStyle]{Lato}
\setmonofont[Mapping=tex-text]{Source Code Pro}

% Large chapter numbers
\usepackage{fix-cm}
\usepackage{titlesec}
\definecolor{gray75}{gray}{0.75}
\newcommand{\hsp}{\hspace{0pt}}
\titleformat{\chapter}[hang]{\flushright\fontseries{b}\fontsize{80}{100}\selectfont}{\fontseries{b}\fontsize{100}{130}\selectfont\textcolor{gray75}{\thechapter}\hsp}{0pt}{\break\Huge\bfseries}[]



%%%%%%%%%%%%%%%%%%%%%%%%%%%%%%%%%%%%%%%%%%%%%%%%%%%%%%%%%%%%%%%%%   Layout-Tuning

\usepackage[bottom]{footmisc}
% For final paragraph spacing/ragged bottom tuning:
%\raggedbottom
%\enlargethispage{\baselineskip}
% + manual linebreaks
% See: http://www.latex-community.org/forum/viewtopic.php?f=47&t=10862


% To enforce placement of floats two options exist, which do not
% produce a misplaced \clearpage. Still the behaviour is different.
% \afterpage delays the float insertion and therefore guarantees
% that the previous page is filled optimally.
% \FloatBarrier will keep a possible change in content better served
% and may leave the previous page partially empty.
\usepackage{placeins}                       % \FloatBarrier
\usepackage{afterpage}                      % \afterpage{\clearpage}

\usepackage{subcaption}
\usepackage[margin=8mm,font=small,labelfont=bf]{caption}        % Indent 8mm, size: small, textbf{figure}
% \renewcommand{\captionfont}{\small\slshape}                   % Font: small italics
% \addtokomafont{caption}{\small}                               % Another way to say: small font



%%%%%%%%%%%%%%%%%%%%%%%%%%%%%%%%%%%%%%%%%%%%%%%%%%%%%%%%%%%%%%%%%   Colors
\definecolor{MidnightBlue}{RGB}{0,0,160} 
\definecolor{blue1}{RGB}{0,0,160} 
\definecolor{green1}{RGB}{0,160,0} 
\definecolor{red1}{RGB}{160,0,0} 
\definecolor{orange1}{RGB}{229,92,0} 
\definecolor{blue120}{RGB}{204,204,236} 
\definecolor{green120}{RGB}{204,236,204} 
\definecolor{red120}{RGB}{236,204,204} 

\definecolor{Crimson}{rgb}{0.6471, 0.1098, 0.1882}
\definecolor{DarkBlue}{rgb}{0.1216, 0.1843, 0.2902}
\definecolor{Silver}{rgb}{0.5490, 0.5451, 0.5882}
\definecolor{darkblue}{rgb}{0.1216, 0.1843, 0.44}


%%%%%%%%%%%%%%%%%%%%%%%%%%%%%%%%%%%%%%%%%%%%%%%%%%%%%%%%%%%%%%%%%   UTF-8 character definitions
%% Compare http://tex.stackexchange.com/questions/29459/declareunicodecharacter-doesnt-work-for-all-characters
%% with respect to DeclareUnicodeCharacter
\usepackage{amsfonts}                   %% math packages
\usepackage{amssymb}
\usepackage{amsmath}

\usepackage{textcomp}
\usepackage{gensymb}
\usepackage{newunicodechar}

\newunicodechar{ℝ}{\mathbb{R}}
\newunicodechar{…}{\ldots}
\newunicodechar{°}{\degree}



%%%%%%%%%%%%%%%%%%%%%%%%%%%%%%%%%%%%%%%%%%%%%%%%%%%%%%%%%%%%%%%%%   Draft stage helpers

% For a really nice todo package including colours & listings, see
% https://tex.stackexchange.com/questions/9796/how-to-add-todo-notes
\usepackage[colorinlistoftodos,prependcaption,textsize=tiny]{todonotes}

\usepackage{ifthen}
\usepackage{comment}                 % Allows make comments out of environments
\usepackage{etoolbox}                % Easy test for empty arguments

%% https://en.wikibooks.org/wiki/LaTeX/Colors
%% https://en.wikibooks.org/wiki/LaTeX/Macros

%% Variable to control output of title page, additional lists, appendices, etc.
%% Could be made a package with DeclareOption which alters the variable.
\newboolean{fulldocument}
\setboolean{fulldocument}{false}

\newcommand{\onlyinfulldocument}[1]{%       % Suppress output if not in internaldraft mode
  \ifthenelse{\boolean{fulldocument}}{%
    #1
  }{}%
}%

%% Could be made a package with DeclareOption which alters the variable.
\newboolean{internaldraft}
\setboolean{internaldraft}{true}

\newcommand{\onlyindraft}[1]{%       % Suppress output if not in internaldraft mode
  \ifthenelse{\boolean{internaldraft}}{%
    #1
  }{}%
}%

\ifthenelse{\boolean{internaldraft}} {%
  \newenvironment{outline}{%
    \textcolor{MidnightBlue}\bgroup
    \textbf{Section Outline:}\\
  }{ \egroup}%

  \newenvironment{itemoutline}{%
    \textcolor{MidnightBlue}\bgroup
    \textbf{Section Outline:}\\
    \begin{itemize}
  } {%
    \end{itemize} \egroup
  }%

} {%
  \excludecomment{outline}
  \excludecomment{itemoutline}
}%

\newcommand{\sidemark}{\onlyindraft{\marginpar{\rule[-3mm]{2mm}{6mm}}}}
\newcommand{\sideremark}[1]{\onlyindraft{\marginpar{\rule[-7mm]{2mm}{10mm}~\parbox[t]{32mm}{#1}}}}
\newcommand{\sidenote}[1]{\onlyindraft{\marginpar{#1}}}


%\newcommand{\REMARK}[1]{\sideremark{#1}}
\newcommand{\REMARK}[1]{\onlyindraft{\todo[linecolor=green,backgroundcolor=green!25,bordercolor=green]{#1}}}

%% Marker where I left writing
\newcommand{\TODO}[1][]{
  \onlyindraft{%
    \textcolor{red}{\ldots}%
    \todo[linecolor=red,backgroundcolor=red!25,bordercolor=red]{%
      TODO: \ifstrempty{#1}{to be continued}{#1}%
    }%
  }%
}%

%% Marker for additional citations
\newcommand{\NEEDSCITATION}[1][]{
  \onlyindraft{%
    \todo[linecolor=blue,backgroundcolor=blue!25,bordercolor=blue]{%
      Citation needed\ifstrempty{#1}{}{: #1}%
    }%
  }%
}%


% %%% Draft revision mark at end of every page %%%%%%%%%%%%
%
% The AddEverypageHook is used to register the printrev command.
% An optional parameter allows arbitrary revision tags.
% \printrev is conditioned on the bool internaldraft.
% When internaldraft is false, the marker is not printed.
\usepackage{everypage}
\usepackage[absolute]{textpos}
\usepackage{datetime}
\newcommand\printrev[1]{%
  \onlyindraft{%
    \begin{textblock*}{\paperwidth}(0mm,\paperheight-16mm)%
      \begin{center}%
        \textit{\footnotesize --- Document rendered on \today{} \currenttime{} #1 --- Build \#\input{counter.dat} ---}%
      \end{center}%
    \end{textblock*}%
  }
}
\AddEverypageHook{\printrev{}}


%%%%%%%%%%%%%%%%%%%%%%%%%%%%%%%%%%%%%%%%%%%%%%%%%%%%%%%%%%%%%%%%%   Title page
% Titlepage and below: courtesy of Jeremie Papon 
% with modifications by Jan-Matthias Braun
% For details of forking and changes, please check the header.

% Trying to unstrip the KOMA title page.
% Although this does not seem to work in any way.
\newcommand*{\mytitlemargin}{\dimexpr 0mm\relax}
\renewcommand{\coverpagetopmargin}{\mytitlemargin}
\renewcommand{\coverpageleftmargin}{\mytitlemargin}
\renewcommand{\coverpagerightmargin}{\mytitlemargin}
\renewcommand{\coverpagebottommargin}{2\mytitlemargin}

% some definitions
\def\title#1{\gdef\title{#1}}
\def\author#1{\gdef\author{#1}}
\def\degreeyear#1{\gdef\degreeyear{#1}}
\def\degreemonth#1{\gdef\degreemonth{#1}}
\def\degree#1{\gdef\degree{#1}}
\def\advisor#1{\gdef\advisor{#1}}
\def\department#1{\gdef\department{#1}}
\def\field#1{\gdef\field{#1}}
\def\university#1{\gdef\university{#1}}
\def\universitycity#1{\gdef\universitycity{#1}}
\def\universitystate#1{\gdef\universitystate{#1}}

\renewcommand{\maketitle}{%
  \begin{titlepage}
  %\singlespacing
  %\thispagestyle{empty}
  \begin{flushright}
  %\vspace*{100pt} 
%  \vspace*{20pt} 
  %\vspace*{\fill}
  \Huge \textcolor{DarkBlue}{\textit{\title}} %\vspace{10pt}
  \textcolor{Silver}{\rule{\textwidth}{2.0pt}}   
  ~ \\ %%%
  \normalsize \normalfont \scshape Dissertation \\ 
  \vspace{6pt}
  in order to obtain the doctoral degree\\
  %in Mathematics and Natural Sciences\\ 
  "Doctor rerum naturalium" \\
  %%% \vspace{6pt}
  of the Georg-August-Universität Göttingen \\
  \vspace{12pt} %%%
  \vspace{6pt}
  in the Doctoral program \\
  Theoretical and Computational Neuroscience (PTCN) of\\
  the Georg-August University School of Science (GAUSS) \\
  \vspace{24pt}
  \vspace{4pt} %%%
  submitted by\\
  \author\\
  \vspace{12pt} 
  of Göttingen, Germany\\
  (place of birth)\\
  \vspace{36pt}
  \vspace*{\fill}
  \includegraphics[width=30mm]{figures/UniLogo1} \\
  \vspace{12pt} 
  \university\\ \universitycity, \universitystate\\
  \degreemonth\ \degreeyear \\
  \end{flushright}
  %\vspace*{\fill}
  \end{titlepage}
}

\newcommand{\abstractpage}{
  \newpage
  %\pagenumbering{roman}
  %\setcounter{page}{3}
  \thispagestyle{empty}
%  \pagestyle{fancy}
%  \renewcommand{\headrulewidth}{0.0pt}
  \vspace*{-28mm}
  \begin{center}
  \vspace*{1pt} % \huge
  \LARGE \textcolor{DarkBlue}{\textit{\title}} \\
  %\vspace*{15pt}
  \vspace*{9pt}
  \Large \scshape Abstract \\ \normalfont \rmfamily
  \end{center}
  \singlespacing

Ramble ramble ramble …

\onehalfspace

%%% Local Variables:
%%% ispell-local-dictionary: "british"
%%% coding: utf-8
%%% mode: latex
%%% TeX-master: "dissertation"
%%% TeX-engine: xetex
%%% End:

  \newpage
  %\cfoot{\thepage}
}


%%%%%%%%%%%%%%%%%%%%%%%%%%%%%%%%%%%%%%%%%%%%%%%%%%%%%%%%%%%%%%%%%   Hyperref at the end
% Funny options by Mr. Papon
% There is supposed to be a setup command which is 
% better in terms of encoding issues.
\RequirePackage[bookmarks, 
                colorlinks=true, 
                citecolor=Crimson,
                filecolor=black,
                linkcolor=DarkBlue,
                urlcolor=Crimson,
                pdfdisplaydoctitle=true,
                ]{hyperref} 


%% Only for T1 text-encoding
%% Problem: There is no direct translation from Unicode to T1
%\DeclareTextCommandDefault{\nobreakspace}{\leavevmode\nobreak\ } 
